%\begin{frame}{Retour sur le théorème CAP}
%
%    Faire graphe et placer les différents systèmes (RDBMS, NoSQL, etc.)
%
%    \note{
%        Utiliser cette transition pour reconnecter la typologie au théorème CAP.
%        Les familles NoSQL n'imposent pas toutes le même compromis :
%        il faut examiner le produit, sa configuration et son architecture.
%        Le diagramme CAP est une simplification pédagogique, pas une étiquette
%        absolue pour tous les déploiements.
%    }
%
%\end{frame}

\begin{frame}{Retour sur le théorème CAP}

    \begin{figure}[htb]
        \resizebox{!}{0.85\textheight}{
        \begin{tikzpicture}
            % Circle 1: Consistency
            \draw[very thick, color=color-consistency, pattern=north east lines, pattern color=color-consistency] (0,0) circle (3cm);
            \node[fill=color-consistency, text=white, font=\bfseries, rounded corners] at (0,0.5) {Consistency};

            % Circle 2: Availability
            \draw[very thick, color=color-availability] (-2.1,-3.25) circle (3cm);
            \node[fill=color-availability, text=white, font=\bfseries, rounded corners] at (-2.75,-3.5) {Availability};

            % Circle 3: Partition Tolerance
            \draw[very thick, color=color-partition] (2.1,-3.25) circle (3cm);
            \node[fill=color-partition, text=white, font=\bfseries, rounded corners, align=left] at (2.75,-3.5) {Partition \\ tolerance};

            \node[fill=none, font=\bfseries, text=color-ca] (ca) at (-1.4,-1.5) {CA};
            \node[fill=none, font=\bfseries, text=color-cp] (cp) at (1.4,-1.5) {CP};
            \node[fill=none, font=\bfseries, text=color-ap] (ap) at (0,-3.7) {AP};

            % Unicorn
            \node[inner sep=0pt] (unicorn) at (0,-2.25) {\includegraphics[width=1cm]{img/unicorn.png}};

            % Relationships
            % AP
            \visible<5->{\node[yshift=-3.1cm] at (ap.south) {\includesvg[width=3cm]{img/db/cassandra.svg}};}
            % CP
            \visible<4->{\node[xshift=3.25cm, yshift=2.5cm] at (cp.east) {\includesvg[width=3cm]{img/db/mongodb.svg}};}
            \visible<3->{\node[xshift=4cm, yshift=1.2cm] at (cp.east) {\includesvg[width=3cm]{img/db/redis.svg}};}
            % CA
            \visible<2->{\node[xshift=-2cm, yshift=3.3cm] at (ca.west) {\includesvg[width=1.6cm]{img/db/postgresql.svg}};}
            \visible<2->{\node[xshift=-3cm, yshift=2cm] at (ca.west) {\includesvg[width=2cm]{img/db/mysql.svg}};}
            \visible<6->{\node[xshift=-4.2cm, yshift=0.5cm] at (ca.west) {\includesvg[width=2cm]{img/db/neo4j.svg}};}
        \end{tikzpicture}
        }
    \end{figure}

    \note{
        Lire le diagramme en rappelant que Cassandra est présenté comme AP,
        MongoDB et Redis comme exemples de systèmes privilégiant la cohérence
        et la tolérance aux partitions, et les SGBD relationnels classiques
        comme exemples CA dans le cadre simplifié du cours.
        Neo4j est placé ici à titre pédagogique : son positionnement réel dépend
        de l'architecture et du mode de déploiement.
        La zone centrale avec la licorne rappelle qu'on ne garantit pas
        simultanément les trois propriétés dans le modèle CAP.
    }
\end{frame}


%\begin{frame}{Performance}
%
%    Faire un tableau comparatif des performances des différents types de bases de données NoSQL avec un exemple.
%
%    Exemple : Quel est le revenu moyen des utilisateurs par pays ?
%
%    SQLite (RDBMS) vs MongoDB (Document) vs Cassandra (Column) vs Neo4j (Graph)
%
%    \note{
%        Faire comparer une même question selon le modèle de données.
%        Pour un revenu moyen par pays, SQLite ou MongoDB peuvent être naturels
%        si les agrégations sont centrales. Cassandra exige une table conçue
%        pour cette requête, et Neo4j est pertinent seulement si les relations
%        entre utilisateurs, pays et transactions sont au cœur du besoin.
%        Il n'existe pas de classement universel : la performance dépend du modèle,
%        des index, du volume et de la requête.
%    }

%\end{frame}

%\begin{frame}
%    Tableau comparatif des 4 types avec :
%
%    Cas d’usage typiques (ex : clé-valeur = cache, graphe = réseaux sociaux).
%    Exemples de bases (Redis, MongoDB, Cassandra, Neo4j).
%    Avantages/Inconvénients (reprise de 3.7_avantages_inconvenients.tex mais spécifique à chaque type).
%
%    Format :
%
%    | Type          | Cas d'usage          | Exemples       | Avantages               | Inconvénients          |
%    |---------------|----------------------|----------------|-------------------------|------------------------|
%    | Clé-valeur    | Cache, session       | Redis, DynamoDB | Rapidité, simplicité   | Pas de requêtes complexes |
%    | Document      | Données semi-structurées | MongoDB, CouchDB | Flexibilité, JSON      | Pas de jointures        |
%    | Colonne       | Big Data, analytique | Cassandra, HBase | Scalabilité, performance | Complexité de modélisation |
%    | Graph         | Réseaux sociaux, recommandations | Neo4j, ArangoDB | Relations complexes, requêtes de graphe | Moins performant pour les données tabulaires |
%
%    \note{
%        Utiliser ce tableau pour faire formuler le compromis de chaque famille.
%        Clé-valeur privilégie l'accès direct, document la flexibilité,
%        colonne la distribution et le débit, graphe les relations.
%        Terminer en rappelant que le cas d'usage et les requêtes doivent guider
%        le choix avant le nom du produit.
%    }
%
%\end{frame}

\begin{frame}{Choisir une famille selon le cas d'usage}
    \small
    \begin{tabular}{p{.17\textwidth}p{.45\textwidth}p{.28\textwidth}}
        \toprule
        \textbf{Famille} & \textbf{Cas d'usage} & \textbf{Exemples} \\
        \midrule
        Clé-valeur & Cache, sessions, classements en temps réel & Redis, DynamoDB \\
        Document & Catalogues, profils, contenus de structure variable & MongoDB, CouchDB \\
        Colonnes & Mesures IoT, journaux, historiques à fort débit d'écriture & Cassandra, HBase \\
        Graphe & Réseaux sociaux, recommandations, chemins, dépendances & Neo4j, Neptune \\
        \bottomrule
    \end{tabular}

    \vspace{0.3cm}

    La forme des données ne suffit pas : selon les requêtes à servir, les mêmes données peuvent justifier plusieurs familles. Le volume, la cohérence attendue, la disponibilité et la latence comptent aussi.

    \note{
        Exemple : des commandes peuvent être stockées dans une base document pour
        afficher une commande complète, dans une base à colonnes pour l'historique
        d'un$\cdot$e client$\cdot$e à grande échelle, ou dans un graphe pour recommander des produits.
        Le critère principal reste le modèle d'accès.
    }
\end{frame}

\begin{frame}{MongoDB ou Cassandra ?}
    \small
    \begin{tabular}{p{.21\textwidth}p{.35\textwidth}p{.35\textwidth}}
        \toprule
         & \textbf{MongoDB} & \textbf{Cassandra} \\
        \midrule
        Architecture & Master-slave & Peer-to-peer \\
        Réplication & Primaire et secondaires, élection en cas de panne & Sans maître, toute réplique accepte les écritures \\
        Théorème CAP & Plutôt CP & Plutôt AP, cohérence ajustable \\
        Modèle & Documents, schéma flexible & Tables conçues pour les requêtes \\
        Requêtes & Variées : index, agrégations & Ciblées par la clé de partition \\
        Cas d'usage & Catalogues, contenus, profils & Écritures massives, séries temporelles \\
        \bottomrule
    \end{tabular}

    \note{
        Pour un catalogue de produits avec des recherches variées, MongoDB est
        généralement préférable, car les requêtes ne sont pas toutes connues à l'avance.
    }
\end{frame}

\begin{frame}{Récapitulatif des architectures}
    \small
    \begin{tabular}{p{.13\textwidth}p{.13\textwidth}p{.17\textwidth}p{.18\textwidth}>{\raggedright\arraybackslash}p{.27\textwidth}}
        \toprule
        \textbf{Base} & \textbf{Famille} & \textbf{Architecture} & \textbf{Écritures} & \textbf{En cas de panne} \\
        \midrule
        Redis & Clé-valeur & Master-slave & Master & Promotion d'un réplica \\
        MongoDB & Document & Master-slave & Primaire & Élection d'un primaire \\
        Cassandra & Colonnes & Peer-to-peer & Toute réplique & Pas de bascule nécessaire \\
        Neo4j & Graphe & Master-slave & Leader & Élection d'un leader \\
        \bottomrule
    \end{tabular}

    \note{
        La dernière colonne décrit la panne du nœud qui reçoit les écritures ;
        chez Cassandra, les autres répliques continuent simplement de répondre.
        Redis Cluster et MongoDB shardé comptent plusieurs masters, un par slot
        pour Redis et un par shard pour MongoDB, mais chaque donnée n'a qu'un
        seul point d'écriture. Chez Neo4j, il existe un leader par base, et le
        cluster est réservé à l'édition Enterprise.
    }
\end{frame}
